% 第九节：转录自“9.12中午_误差分析完成.docx”，统一为现有LaTeX格式。
% 原表7保留数据；原式48、原图9改由LaTeX自动编号及引用。
为评估计算结果的可靠性，分别检验数值实现、离散误差及关键参数扰动的影响。
检验沿用原计算数据，重点考察温度、含水率和临界烘干时长。

\subsection{数值检验与误差分析}
\label{sec:verification-errors}

采用独立对照与网格细化检验数值精度：
问题一温度场与Bessel--Duhamel半解析解对照；
非线性传质方程通过制造解检验\cite{mms-verification}；
各问进一步比较网格及时间步长调整后的结果。
主要指标汇总于表\ref{tab:verification}。

\begin{table}[!htbp]
  \centering
  \caption{数值检验的主要结果}
  \label{tab:verification}
  \begin{tabularx}{\textwidth}{@{}C{0.19\textwidth}YC{0.25\textwidth}@{}}
    \toprule
    检验对象 & 关键结果 & 说明\\
    \midrule
    问题一温度场
      & 最大偏差$1.68\times10^{-6}\tempunit$
      & 半解析解对照\\
    非线性传质
      & 误差$2.088\times10^{-6}\cunit$；阶约$2.00$
      & 制造解检验\\
    问题二含水率
      & 离散误差估计约$2.5\times10^{-5}\cunit$
      & 空间与时间误差之和\\
    问题四含水率
      & 最细误差估计约$1.0\times10^{-4}\cunit$
      & 固定物理位置重采样\\
    长期临界时长
      & 问题三$57.4721\ \mathrm{h}$；问题四$51.0905\ \mathrm{h}$
      & 882节点结果\\
    \bottomrule
  \end{tabularx}
\end{table}

问题一温度与含水率场均呈二阶收敛。
问题三由442细化至882节点，临界时长变化约$0.0003\ \mathrm{h}$；
问题四在882节点上将步长加倍，时长变化约$0.0001\ \mathrm{h}$，
说明达标时间对所检验的离散设置较稳定。
问题四重采样场的观测阶约为$1.63$，
故部分含水率结果的第四位小数仍有不确定性。

离散收支检验中，问题一至四的传质相对残差均不超过$1.78\times10^{-12}$，
问题一、二的热量收支相对残差均不超过$1.63\times10^{-10}$。
该检验对应本文的离散显热模型；收缩域按材料坐标的逐步收支计算。
数值场范围、严格达标条件及正文结果表与工作簿亦已核对。

现有数据缺少内部温度、径向含水率及独立烘干终点观测，
因此上述结果验证的是数值实现与离散精度，
尚不能代表相对实测值的预测误差。

\subsection{参数灵敏度分析}
\label{sec:sensitivity}

以临界烘干时长$t_*$为响应量，
对换热系数$h$、传质系数$h_m$和扩散前置系数$D_0$
分别施加$\pm10\%$的单因素扰动，其余参数保持不变。
归一化灵敏度定义为
\begin{equation}
  S_p^{\pm}
    =\frac{t_*\!\left(p(1\pm0.1)\right)-t_*(p)}
      {\pm0.1\,t_*(p)},
  \qquad p\in\{h,h_m,D_0\}.
  \label{eq:parameter-sensitivity}
\end{equation}

两问各计算1组基准和6组扰动工况，均采用222节点及原时间步长。
基准时长分别为$57.4706\ \mathrm{h}$和$51.0906\ \mathrm{h}$，
与882节点结果相差小于$0.002\ \mathrm{h}$。
$D_0$扰动仅改变前置系数；
问题四保持给定半径历程$R(t)$不变。
各参数引起的时长相对变化见图\ref{fig:parameter-sensitivity}。

\pfig{图12.png}{0.98}{参数增减10\%时临界烘干时长的相对变化}{fig:parameter-sensitivity}

两问均表现为扩散前置系数影响最大、传质系数次之、换热系数最小。
$D_0$增加$10\%$时，烘干时长减少约$8\%$；
降低$10\%$时，时长增加约$9.8\%$。
传质系数引起的时长变化约为$1.0\%$--$1.3\%$（问题三）
和$0.4\%$--$0.6\%$（问题四），
换热系数仅约$0.01\%$--$0.02\%$。

扩散参数扰动的影响远大于网格细化误差，
说明后续提高预测精度应优先标定扩散参数，
尤其是低含水率阶段的扩散能力。
该结论限于题给工况及给定收缩历程，
不能直接用于确定最优温度或风速。
